\section*{Chapitre \ref{ch:ml:build-an-order-book-model-end-to-end} --- Construction~: un modèle de carnet d'ordres, de bout en bout}

\begin{solution}{exo:ml:build-an-order-book-model-end-to-end:1}
Parce que la production voit le marché à travers le flux propre de l'agent, avec sa latence, et calcule ses variables à partir de ce qu'elle a
reçu~; les données de recherche enregistrées par le même agent avec le même code de variables sont celles que verra la production, et les valeurs hors ligne et
en ligne peuvent être testées égales. Un tape complet vu de nulle part entraînerait le modèle sur une information, et un timing, que l'agent
ne possède jamais.
\end{solution}

\begin{solution}{exo:ml:build-an-order-book-model-end-to-end:2}
Acheter au cours vendeur et vendre au cours acheteur une seconde plus tard coûte la fourchette, soit un tick, auquel s'ajoutent deux fois les frais du preneur, $2\times\$0.003=\$0.006$
par action, soit 0.6 tick à un centime~: 1.6 tick par aller-retour. Le mid bouge d'un tick ou plus en une seconde dans environ 5~\% des cas, et
les prévisions sont bien plus petites qu'un tick~; aucun seuil ne laisse de transactions qui couvrent le coût.
\end{solution}

\begin{solution}{exo:ml:build-an-order-book-model-end-to-end:3}
Sans modèle, les cotations sont exécutées 58.7 fois par session avec un mark-out à une seconde de 0.078 tick~: la plupart des exécutions proviennent d'ordres
qui font bouger le prix contre la cotation. Le modèle retire le côté sur le point d'être touché, de sorte que la cotation est exécutée moins souvent (34.3 fois) mais
ses exécutions valent 0.405 tick une seconde plus tard. P\&L par session~: 21.73~\$ contre 5.47~\$.
\end{solution}

\begin{solution}{exo:ml:build-an-order-book-model-end-to-end:4}
Les décisions sont déclenchées à chaque changement du carnet d'ordres, environ quatre par seconde en moyenne mais bien davantage par rafales. Avec $L=100$ millisecondes,
$\lambda L$ franchit un dans les rafales et les décisions attendent les unes derrière les autres, si bien que leur péremption médiane (154 millisecondes) dépasse
l'instant de décision. Un système de production regroupe les mises à jour~: quand il achève une décision, il calcule la suivante à partir du carnet le plus récent, en sautant
les mises à jour qu'il a manquées, si bien que la péremption est bornée par environ deux fois l'instant de décision~; mieux encore, il rend la décision assez rapide
pour que $\lambda L$ reste bien en dessous de un dans les rafales les plus intenses.
\end{solution}

\begin{solution}{exo:ml:build-an-order-book-model-end-to-end:5}
Non~: le modèle a été rentable dans les six sessions, et les grandes valeurs du moniteur viennent des variables de comptage, qui
suivent l'activité de la session (2.8 à 6.4 changements du carnet par seconde), tandis que les variables de prix restent entre 0.02 et 0.13. Le
moniteur devrait utiliser une référence couvrant l'éventail de l'activité (davantage de sessions d'entraînement, ou un historique), normaliser les variables de comptage par
l'activité de la session, ou surveiller les comptages avec un moniteur d'activité distinct qui alerte le desk, et non le \omterm{def:ml:the-machine-learning-team:roles}{propriétaire du modèle}.
\end{solution}

\begin{solution}{exo:ml:build-an-order-book-model-end-to-end:6}
La latence ne vaut de l'argent que là où les décisions sont en concurrence avec les événements ou d'autres traders. Ici le P\&L est identique à 150 nanosecondes
et à 263 microsecondes (21.73~\$), et inchangé à l'intérieur de son erreur type jusqu'à 50 millisecondes~: l'accélération n'apporte rien dans
ce marché. Elle achète de la marge (un modèle qui reste rapide dans les rafales) et elle achète tout dans un marché où les concurrents
rivalisent à la microseconde.
\end{solution}

\begin{solution}{exo:ml:build-an-order-book-model-end-to-end:7}
Les sessions de test modifient leur meilleur niveau du carnet environ 4.3 fois par seconde, et le P\&L franchit zéro vers 106 millisecondes, soit
0.45 changement du carnet. À 2\,000 changements par seconde, le même nombre de changements prend $106\times4.3/2000\approx0.23$ millisecondes, soit environ
230 microsecondes. L'estimation laisse de côté la concurrence (d'autres traders rapides saisissent d'abord les cotations périmées, ce qui déplace la falaise jusqu'à
l'écart entre leur vitesse et celle de l'agent), la position dans la file (les marchés plus rapides ont des files plus longues), le modèle différent
qu'un marché plus actif exigerait, et le caractère en rafales qui fait grossir la file de décisions.
\end{solution}

\begin{solution}{exo:ml:build-an-order-book-model-end-to-end:8}
Entrées~: les quatre variables de prix face à une référence couvrant de nombreuses sessions, avec un seuil calibré pour une fausse alerte par mois
(chapitre~27)~; les variables de comptage face à l'activité de la session, sous la responsabilité du desk. Sorties~: la part des décisions avec un côté
retiré, le taux de disponibilité des cotations, la distribution de la péremption des décisions par rapport à un budget (médiane et 99\up{e}~centile, avec une alarme lorsque
$\lambda L$ dans les rafales s'approche de un). Résultats~: le mark-out à une seconde des exécutions et le P\&L de la session, avec un CUSUM du
mark-out quotidien, sous la responsabilité du \omterm{def:ml:the-machine-learning-team:roles}{propriétaire du modèle}. Service~: parité du chemin C++ avec la référence sur un échantillon quotidien d'entrées en production, sous la responsabilité de
l'\omterm{def:ml:the-machine-learning-team:roles}{ingénieur en apprentissage automatique} (ML engineer).
\end{solution}

\begin{solution}{pb:ml:build-an-order-book-model-end-to-end:1}
\textbf{Partie I.}
\begin{enumerate}
\item Une plateforme du simulateur du livre~10, une action à 100~\$ avec un tick d'un centime, un rabais de 0.002~\$ et des frais de preneur de 0.003~\$ par action,
flux de fond \texttt{firm\_tape}, des latences de 20 microsecondes~; à chaque changement du carnet d'ordres, dix variables et une prévision à une seconde~; coter un lot
à chaque meilleur cours autorisé, en retirant le côté que la prévision annonce comme devant être touché ($\theta=0.1$ ticks), en clôturant les positions après 30
secondes~; huit sessions d'entraînement, trois de validation et six de test, de vingt minutes chacune.
\item Déséquilibre au meilleur cours, fourchette, déséquilibre du flux d'ordres (OFI) sur 1 et 5 secondes, volume signé sur 1 et 5 secondes, transactions sur 5
secondes~; variation du mid sur 1 et 5 secondes, mises à jour sur 1 seconde~; hors ligne et en ligne concordent parce que le même agent et les mêmes
\omterm{def:ml:data-and-feature-stores:store}{définitions de variables} calculent les deux.
\item La variation du mid en ticks sur la seconde suivante, à partir du dernier événement connu à l'instant de décision~; le mid ne bouge généralement pas en une
seconde (89.4~\% de zéros).
\item Validation croisée purgée à cinq plis avec un embargo d'une seconde~; 7 feuilles (0.558 contre 0.555 et 0.548).
\end{enumerate}
\textbf{Partie II.}
\begin{enumerate}[resume]
\item Coefficient d'information de validation (IC) de 0.522 pour la forêt, 0.452 pour le TCN.
\item 153 nanosecondes à la médiane pour les branches générées, contre 263 microsecondes pour l'appel Python de LightGBM et 166 pour
le TCN.
\item Les tableaux plats, la boucle C++ et les branches ainsi que le noyau Rust reproduisent exactement les prédictions de LightGBM sur 200 vecteurs partagés.
\item Sur les sessions de validation, à la latence du modèle compilé~: 53.9~\$ par session à 0.1 tick, le meilleur des quatre.
\end{enumerate}
\textbf{Partie III.}
\begin{enumerate}[resume]
\item 21.73~\$ par session (erreur type de 3.90~\$), positif dans chacune des six~; la cotation sans modèle rapporte 5.47~\$.
\item Les exécutions retenues par le modèle valent 0.405 tick une seconde plus tard, contre 0.078 sans modèle~: le modèle refuse les exécutions qui
perdent.
\item Les plus grands indices vont de 0.14 à 1.60, tous issus de variables de comptage qui suivent l'activité de la session~; les variables de prix atteignent au plus 0.13.
\item Le délai entre l'événement de marché et le départ de l'ordre~: l'instant de décision, auquel s'ajoute l'attente derrière les décisions antérieures lorsque
elles arrivent plus vite qu'elles ne peuvent être prises.
\end{enumerate}
\textbf{Partie IV.}
\begin{enumerate}[resume]
\item \emph{Les microsecondes, c'est de l'argent.} Par session de vingt minutes, 21.73~\$ de zéro à 263 microsecondes, 21.23~\$ à 1 milliseconde,
21.72~\$ à 10, 20.63~\$ à 30, 20.52~\$ à 50, 12.40~\$ à 70, 1.60~\$ à 100 et $-\$12.27$ à 150~; le mark-out à une seconde passe de
0.405 ticks à 0.148~; l'edge tombe à zéro vers 106 millisecondes, et sous la référence sans modèle vers 89.
\item Parce que le carnet d'ordres change environ quatre fois par seconde et que rien d'autre ne se dispute les cotations~: un retrait avec 50 millisecondes de retard
est rarement trop tardif.
\item La péremption croît de la durée propre de la décision et, dans les rafales, de la file de décisions ($\lambda L>1$), exactement quand le marché
est le plus actif.
\item Beaucoup plus de changements du carnet par seconde et des concurrents rapides~: la falaise se situerait plus près de zéro dans le rapport des taux d'événements, et
encore plus près là où les concurrents rivalisent de vitesse.
\item Pas pour le bien de ce marché~; mais pour la robustesse dans les rafales, et pour les marchés où elle compte.
\item Agents rapides concurrents, un flux de fond réactif (\texttt{firm.agentmkt}), des taux d'événements réalistes et le regroupement des mises à jour.
\item Comme au chapitre~25 (les clés d'étape du graphe comme version du code, la forêt comme artefact) et au chapitre~27 (la spécification
de l'\cref{exo:ml:build-an-order-book-model-end-to-end:8}).
\item Ce que révèlent le taux d'événements du marché et ses concurrents~: ici rien, ailleurs tout.
\end{enumerate}
\end{solution}

\begin{solution}{iq:ml:build-an-order-book-model-end-to-end:1}
Déséquilibre du flux d'ordres au premier niveau du carnet sur de courtes fenêtres, déséquilibre des tailles au premier niveau, la fourchette, le volume échangé signé, les variations récentes du prix milieu
et l'activité, construites comme des quantités stationnaires sur des fenêtres en temps d'événements ou en temps d'horloge, comme le font Cont, Kukanov et Stoikov
ainsi que Kolm, Turiel et Westray.
\iqlookfor{flux d'ordres et déséquilibre, stationnarité.}
\end{solution}

\begin{solution}{iq:ml:build-an-order-book-model-end-to-end:2}
Purger des plis d'entraînement toute ligne dont l'intervalle couvert par l'étiquette chevauche le pli de test, ajouter un embargo après celui-ci, et garder les sessions ou les jours de rôles différents sur des données disjointes~;
choisir les seuils sur des données qui n'ont pas servi à ajuster le modèle.
\iqlookfor{purge, embargo et données de test disjointes.}
\end{solution}

\begin{solution}{iq:ml:build-an-order-book-model-end-to-end:3}
Une seule \omterm{def:ml:data-and-feature-stores:store}{définition de variable}, exécutée par les deux chemins (un \omterm{def:ml:data-and-feature-stores:store}{feature store}), des données de recherche enregistrées par le chemin de code de production, et un
test que les valeurs hors ligne égalent les valeurs en ligne aux mêmes instants de décision, exécuté en intégration continue et sur des échantillons en production.
\iqlookfor{code partagé et test de parité.}
\end{solution}

\begin{solution}{iq:ml:build-an-order-book-model-end-to-end:4}
Sur les exécutions, il recule~: moins d'exécutions, mais de meilleurs mark-outs. Mesurer la courbe de mark-out des exécutions avec et sans le modèle, ainsi que le
P\&L par session, sur les mêmes sessions.
\iqlookfor{sélection adverse et mark-outs.}
\end{solution}

\begin{solution}{iq:ml:build-an-order-book-model-end-to-end:5}
Exporter le modèle sous une forme plate (tableaux d'arbres, poids quantifiés), générer ou écrire un noyau natif sans allocation sur le hot path, et
vérifier la parité bit à bit ou à une tolérance près sur des vecteurs de test communs, et mesurer la latence médiane et la queue de distribution sur la machine cible.
\iqlookfor{compilation, vecteurs de parité, queues mesurées.}
\end{solution}

\begin{solution}{iq:ml:build-an-order-book-model-end-to-end:6}
Trader la stratégie sur une grille d'instants de décision, dans un simulateur aux taux d'événements réalistes et avec des concurrents, repérer où son P\&L
fléchit et où il croise zéro, et consacrer des moyens à la réduction de la latence jusqu'à ce que celle de la stratégie se situe confortablement sur la partie plate, y compris en rafales.
\iqlookfor{une courbe de P\&L mesurée en fonction de la latence.}
\end{solution}
